\chapter{Modelagem UML e modelos de apoio}\label{cap:uml}

\section{Leitura e finalidade dos modelos}
Os modelos representam perspectivas diferentes do mesmo domínio e usam a
notação UML como referência \cite{uml251}. Casos de uso mostram objetivos dos
atores; classes mostram conceitos; sequência mostra mensagens; atividades
mostram decisões; estados mostram o ciclo de vida de uma inscrição.
Os contratos e as fichas continuam definindo as obrigações do produto.

\begin{quadrosrs}
\caption{Perguntas respondidas pelos modelos}
\begin{tabelasrs}{L{0.25\textwidth}Y}
\cabecalhosrs\textbf{Modelo} & \textbf{Pergunta e vínculo com a especificação}\\\midrule
Casos de uso & Quem consulta, inscreve e cancela? \ucref{UC-001} a \ucref{UC-004}.\\
Classes & Quais conceitos, atributos e cardinalidades existem? Dicionário da \autoref{tab:dados}.\\
Sequência e atividades & Como a inscrição responde à repetição, recusa e falha? \reqref{RF-002}, \reqref{RF-003} e contrato HTTP.\\
Estados & Como uma inscrição passa de ativa a cancelada? \autoref{fig:uml-estados}, no \autoref{ap:glossario}.\\
Componentes, implantação e pacotes & Como distribuir responsabilidades, execução e dependências? Exemplos de design compatíveis com o domínio.\\
Dados e rastreabilidade & Como relacionar entidades e ligar obrigação a teste? Modelos complementares; DER e grafo de rastreabilidade não são diagramas UML.\\
\end{tabelasrs}
\fonte{Elaboração didática deste guia.}
\end{quadrosrs}

Uma linha de associação liga ator e objetivo ou conceitos relacionados.
Multiplicidades indicam quantas instâncias participam de uma relação.
Setas tracejadas representam dependência nos modelos de design e retorno
no diagrama de sequência. Guardas entre colchetes delimitam caminhos.
O traço \texttt{/} antes de um atributo indica valor derivado.

\clearpage
\section{Casos de uso e fronteira}
A \autoref{fig:uml-casos} reúne os quatro casos já especificados no
\autoref{cap:casos}. Os atores ficam fora da fronteira do sistema.
Participante autenticado também pode consultar o catálogo público.
Autenticação é condição das operações pessoais; não é um cadastro de contas
implementado pelo produto. Duplicidade é uma alternativa de \ucref{UC-001},
conforme \reqref{RF-003}, e não um objetivo independente do ator.

\diagramasrs{UML de casos de uso do sistema de inscrições}{fig:uml-casos}
  {\begin{tikzpicture}
  \draw[draw=srslinha] (1.8,-5.0) rectangle (8.0,4.5);
  \node[font=\small\bfseries] at (4.9,4.15) {Sistema de inscrições};
  \atoruml{visitante}{(0,2.8)}{Visitante}
  \atoruml{participante}{(0,-1.6)}{Participante\\autenticado}
  \node[uml/caso] (catalogo) at (4.9,2.8) {\ucref{UC-003}\\Consultar eventos};
  \node[uml/caso] (criar) at (4.9,0.7) {\ucref{UC-001}\\Realizar inscrição};
  \node[uml/caso] (consultar) at (4.9,-1.5) {\ucref{UC-004}\\Consultar próprias\\inscrições};
  \node[uml/caso] (cancelar) at (4.9,-3.8) {\ucref{UC-002}\\Cancelar inscrição};
  \draw (visitante.east) -- (catalogo.west);
  \draw (participante.east) -- (criar.west);
  \draw (participante.east) -- (consultar.west);
  \draw (participante.east) -- (1.5,-1.6) -- (1.5,-3.8) -- (cancelar.west);
  \draw (participante.north east) -- (1.1,-0.3) -- (1.1,1.9) -- (catalogo.south west);
  \node[uml/nota,text width=39mm] at (10.5,1.7)
    {A consulta de eventos é pública, inclusive para quem já possui sessão.};
  \node[uml/nota,text width=39mm] at (10.5,-1.8)
    {Catálogo e identidade são integrações externas; sua administração fica fora dos casos de uso.};
\end{tikzpicture}
}

\clearpage
\section{Classes do domínio}
A \autoref{fig:uml-classes} representa conceitos e operações do domínio,
sem controllers, tabelas de senhas ou classes de framework.
Participante representa a identidade externa, não uma obrigação de armazenar
perfil local. Cada inscrição tem um evento e um titular; cada evento ou titular
pode ter vários registros históricos, inclusive cancelados.

\diagramasrs{UML de classes conceituais e cardinalidades}{fig:uml-classes}
  {\begin{tikzpicture}
  \node[uml/classe,rectangle split parts=2] (pessoa) at (0,0)
    {\textbf{Participante}
    \nodepart{second}+ id: Identificador};
  \node[uml/classe] (evento) at (8.6,0)
    {\textbf{Evento}
    \nodepart{second}+ id: Identificador\\+ nome: Texto\\+ inicio: InstanteUTC\\
      + capacidade: Inteiro\\+ publicado: Booleano\\+ /vagas: Inteiro
    \nodepart{third}+ elegivel(agora): Booleano};
  \node[uml/classe,text width=57mm] (inscricao) at (3.6,-6.0)
    {\textbf{Inscrição}
    \nodepart{second}+ id: Identificador\\+ eventoId: Identificador\\
      + participanteId: Identificador\\+ situacao: SituacaoInscricao\\
      + criadaEm: InstanteUTC\\+ canceladaEm:\\\hspace*{1em}InstanteUTC [0..1]
    \nodepart{third}+ cancelar(agora)\\+ pertenceA(titular): Booleano};
  \node[uml/caixa,rectangle split,rectangle split parts=2,
    rectangle split part align={center,left},text width=37mm] (situacao) at (9.1,-6.0)
    {\guillemotleft enumeration\guillemotright\\\textbf{SituacaoInscricao}
    \nodepart{second}ativa\\cancelada};
  \draw (pessoa.south) -- node[uml/rotulo,pos=0.10,left]{1}
    node[uml/rotulo,pos=0.88,left]{0..*}
    node[uml/rotulo,midway,left=3mm]{titular} (inscricao.north west);
  \draw (evento.south) -- node[uml/rotulo,pos=0.10,right]{1}
    node[uml/rotulo,pos=0.90,right]{0..*}
    node[uml/rotulo,midway,right=3mm]{evento} (inscricao.north east);
  \draw[uml/dependencia] (inscricao.east) -- (situacao.west);
\end{tikzpicture}
}

A restrição de \rnref{RN-001} limita os registros \textbf{ativos}, não todo o
histórico. \texttt{/vagas} é capacidade menos inscrições ativas.
O atributo \texttt{canceladaEm} tem multiplicidade \texttt{0..1}, respeitando
o dicionário. As operações indicam responsabilidades conceituais; a criação
atômica é responsabilidade do sistema e não de uma coleção manipulada pela UI.

\clearpage
\section{Sequência da inscrição}
A interação da \autoref{fig:uml-sequencia} detalha uma solicitação com corpo
e sessão válidos. Validação recusada interrompe o fluxo antes da leitura pessoal,
como mostram o contrato e a atividade da \autoref{fig:uml-atividades}.
As quatro linhas de vida são papéis lógicos; não estabelecem quatro serviços
implantados. Mensagens contínuas têm ponta preenchida e retornos são tracejados.

\diagramasrs{UML de sequência de UC-001 e seus resultados}{fig:uml-sequencia}
  {\begin{tikzpicture}[font=\footnotesize]
  \atoruml{pessoa}{(0,0)}{Participante}
  \node[uml/caixa,font=\footnotesize,text width=23mm] (ui) at (3.5,0) {:Interface};
  \node[uml/caixa,font=\footnotesize,text width=27mm] (sistema) at (7.1,0) {:Sistema};
  \node[uml/caixa,font=\footnotesize,text width=31mm] (dados) at (11.2,0) {:Armazenamento};
  \foreach \x in {0,3.5,7.1,11.2}
    \draw[dashed,draw=srslinha] (\x,-1.15) -- (\x,-15.9);
  \draw[uml/chamada] (0,-1.6) -- node[uml/rotulo,above]{confirmar evento} (3.5,-1.6);
  \draw[uml/chamada] (3.5,-2.6) -- node[uml/rotulo,above]{POST /inscricoes} (7.1,-2.6);
  \draw[uml/chamada] (7.1,-3.2) -- (8.1,-3.2) -- (8.1,-4.2)
    -- node[uml/rotulo,right,text width=30mm]{validar corpo\\e sessão} (7.1,-4.2);
  \draw[draw=srslinha] (6.4,-4.8) rectangle (13.1,-7.5);
  \node[uml/rotulo,anchor=north west,fill=white] at (6.4,-4.8) {critical};
  \draw[uml/chamada] (7.1,-5.8) -- node[uml/rotulo,above,text width=35mm]
    {obter existente ou\\criar se elegível} (11.2,-5.8);
  \draw[uml/retorno] (11.2,-6.8) -- node[uml/rotulo,above]{resultado da operação} (7.1,-6.8);
  \draw[draw=srslinha] (2.0,-8.0) rectangle (13.1,-14.7);
  \node[uml/rotulo,anchor=north west,fill=white] at (2.0,-8.0) {alt};
  \node[uml/rotulo,anchor=north west] at (3.0,-8.0) {[inscrição ativa existente]};
  \draw[uml/retorno] (7.1,-9.0) -- node[uml/rotulo,above]{200 + mesmo ID} (3.5,-9.0);
  \draw[dashed,draw=srslinha] (2.0,-9.6) -- (13.1,-9.6);
  \node[uml/rotulo,anchor=north west] at (3.0,-9.6) {[nova inscrição confirmada]};
  \draw[uml/retorno] (7.1,-10.6) -- node[uml/rotulo,above]{201 + novo ID} (3.5,-10.6);
  \draw[dashed,draw=srslinha] (2.0,-11.2) -- (13.1,-11.2);
  \node[uml/rotulo,anchor=north west] at (3.0,-11.2) {[operação recusada]};
  \draw[uml/retorno] (7.1,-12.2) -- node[uml/rotulo,above]{404 ou 409; sem gravação} (3.5,-12.2);
  \draw[dashed,draw=srslinha] (2.0,-12.8) -- (13.1,-12.8);
  \node[uml/rotulo,anchor=north west] at (3.0,-12.8) {[falha de leitura ou gravação]};
  \draw[uml/retorno] (7.1,-13.8) -- node[uml/rotulo,above]{503; sem confirmação} (3.5,-13.8);
  \draw[uml/retorno] (3.5,-15.5) -- node[uml/rotulo,above]{apresentar resultado} (0,-15.5);
\end{tikzpicture}
}

O fragmento \texttt{critical} representa o trecho protegido contra concorrência.
O fragmento \texttt{alt} separa resultados alternativos.
Se a resposta se perder após a confirmação, a repetição percorre o caminho
de registro existente; não consome outra vaga, conforme \reqref{RF-003}.

\clearpage
\section{Atividades, decisões e exceções}
O fluxo da \autoref{fig:uml-atividades} liga as guardas de sessão, duplicidade
e elegibilidade. Os nós finais encerram caminhos alternativos da mesma atividade.
Corpo inválido ou sessão inválida impede acesso pessoal; ausência de vaga,
publicação ou prazo impede nova criação. Uma inscrição ativa existente é
consultada antes dessas condições de criação.

\diagramasrs{UML de atividades para criação e repetição da inscrição}{fig:uml-atividades}
  {\begin{tikzpicture}
  \node[circle,fill=black,inner sep=2pt] (inicio) at (0,0) {};
  \node[uml/acao] (validar) at (0,-1.2) {Validar corpo e sessão};
  \node[uml/decisao] (sessao) at (0,-2.7) {Válidos?};
  \node[uml/acao,text width=30mm] (erroentrada) at (-5,-2.7) {Recusar:\\400, 401 ou 422};
  \node[uml/acao] (buscar) at (0,-4.2) {Buscar inscrição ativa\\do titular/evento};
  \node[uml/decisao] (existente) at (0,-5.7) {Existe?};
  \node[uml/acao,text width=30mm] (repetir) at (5,-5.7) {Retornar 200\\com o mesmo ID};
  \node[uml/decisao] (elegivel) at (0,-7.4) {Elegível?};
  \node[uml/acao,text width=30mm] (recusar) at (-5,-7.4) {Recusar:\\404 ou 409};
  \node[uml/acao] (gravar) at (0,-9.0) {Persistir criação e\\efeito na capacidade};
  \node[uml/decisao] (confirmada) at (0,-10.6) {Confirmada?};
  \node[uml/acao,text width=30mm] (indisponivel) at (-5,-10.6) {Informar 503\\sem sucesso};
  \node[uml/acao] (sucesso) at (0,-12.2) {Retornar 201\\com novo ID};
  \finaluml{fim}{(0,-13.6)}
  \finaluml{fimEntrada}{(-5,-4.0)}
  \finaluml{fimRecusa}{(-5,-8.7)}
  \finaluml{fimFalha}{(-5,-12.0)}
  \finaluml{fimRepeticao}{(5,-7.0)}
  \draw[uml/fluxo] (inicio) -- (validar);
  \draw[uml/fluxo] (validar) -- (sessao);
  \draw[uml/fluxo] (sessao) -- node[uml/rotulo,above]{[não]} (erroentrada);
  \draw[uml/fluxo] (sessao) -- node[uml/rotulo,right]{[sim]} (buscar);
  \draw[uml/fluxo] (buscar) -- (existente);
  \draw[uml/fluxo] (existente) -- node[uml/rotulo,above]{[sim]} (repetir);
  \draw[uml/fluxo] (existente) -- node[uml/rotulo,right]{[não]} (elegivel);
  \draw[uml/fluxo] (elegivel) -- node[uml/rotulo,above]{[não]} (recusar);
  \draw[uml/fluxo] (elegivel) -- node[uml/rotulo,right]{[sim]} (gravar);
  \draw[uml/fluxo] (gravar) -- (confirmada);
  \draw[uml/fluxo] (confirmada) -- node[uml/rotulo,above]{[não]} (indisponivel);
  \draw[uml/fluxo] (confirmada) -- node[uml/rotulo,right]{[sim]} (sucesso);
  \draw[uml/fluxo] (sucesso) -- (fim);
  \draw[uml/fluxo] (erroentrada) -- (fimEntrada);
  \draw[uml/fluxo] (recusar) -- (fimRecusa);
  \draw[uml/fluxo] (indisponivel) -- (fimFalha);
  \draw[uml/fluxo] (repetir) -- (fimRepeticao);
  \node[uml/nota,text width=39mm] at (4.9,-2.6)
    {Trecho de domínio: leitura consistente, validação de vaga e gravação na mesma operação atômica.};
  \node[uml/nota,text width=39mm] at (4.9,-9.1)
    {Elegibilidade: evento conhecido, publicado, no prazo e com vaga. Duplicidade é decidida antes do prazo.};
\end{tikzpicture}
}

Leitura consistente, validação e gravação do trecho de domínio fazem parte da
operação atômica de \rnref{RN-001}. Um erro não altera registros nem vagas.
Falha de leitura também encerra a operação com 503, antes de avaliar o resultado.
Os limites exatos de prazo são exercitados por \ctref{CT-002} e \ctref{CT-004}.

\clearpage
\section{Componentes e responsabilidades}
As três vistas de design a seguir são exemplos de organização, independentes
de tecnologia. Não acrescentam infraestrutura obrigatória aos requisitos.
Na \autoref{fig:uml-componentes}, a aplicação coordena os casos de uso,
o domínio reúne regras e a persistência conserva eventos e inscrições.
Uma dependência aponta do consumidor para aquilo que ele usa.

\diagramasrs{UML de componentes em uma decomposição didática}{fig:uml-componentes}
  {\begin{tikzpicture}
  \draw[draw=srslinha,dashed] (-2.2,1.7) rectangle (9.6,-5.5);
  \node[font=\small\bfseries,anchor=west] at (-2.0,1.35) {Responsabilidade do produto};
  \node[uml/componente] (ui) at (0,0)
    {\guillemotleft component\guillemotright\\Interface Web};
  \node[uml/componente] (app) at (5.1,0)
    {\guillemotleft component\guillemotright\\Aplicação};
  \node[uml/componente] (dominio) at (0,-3.9)
    {\guillemotleft component\guillemotright\\Regras do domínio};
  \node[uml/componente] (persistencia) at (5.1,-3.9)
    {\guillemotleft component\guillemotright\\Persistência};
  \node[uml/componente,text width=30mm] (catalogo) at (12.0,0)
    {\guillemotleft component\guillemotright\\Catálogo externo};
  \node[uml/componente,text width=30mm] (identidade) at (12.0,-3.9)
    {\guillemotleft component\guillemotright\\Identidade externa};
  \draw[uml/dependencia] (ui) -- node[uml/rotulo,above,text width=15mm]{HTTP/\\JSON} (app);
  \draw[uml/dependencia] (app.south west) -- node[uml/rotulo,above,sloped]
    {aplica regras} (dominio.north east);
  \draw[uml/dependencia] (app) -- node[uml/rotulo,right]{grava/consulta} (persistencia);
  \draw[uml/dependencia] (app) -- node[uml/rotulo,above]{INT-001} (catalogo);
  \draw[uml/dependencia] (app.south east) -- node[uml/rotulo,above,sloped]
    {INT-002} (identidade.north west);
\end{tikzpicture}
}

Catálogo e identidade continuam fora da responsabilidade do produto.
Se a implementação tiver componentes diferentes, conserve o vínculo com os
contratos INT-001 e INT-002 e com \reqref{RNF-004}.
Separação de componentes não significa, por si só, microserviços ou processos
independentes. As interfaces oferecidas e consumidas podem ser detalhadas em
uma especificação de design associada aos IDs dos requisitos.

\clearpage
\section{Implantação e artefatos}
A \autoref{fig:uml-implantacao} distingue nós de execução, artefatos e caminhos
de comunicação. O serviço e seu armazenamento têm responsabilidades diferentes:
reiniciar o processo da aplicação não remove os registros confirmados.

\diagramasrs{UML de implantação em um ambiente de referência}{fig:uml-implantacao}
  {\begin{tikzpicture}
  \nouml{cliente}{(0,0)}{\guillemotleft device\guillemotright\\Dispositivo cliente}
    {\guillemotleft artifact\guillemotright\\Interface no navegador}
  \nouml{servidor}{(7.0,0)}{\guillemotleft executionEnvironment\guillemotright\\Ambiente da aplicação}
    {\guillemotleft artifact\guillemotright\\Serviço de inscrições}
  \nouml{armazenamento}{(7.0,-5.0)}{\guillemotleft node\guillemotright\\Armazenamento durável}
    {\guillemotleft artifact\guillemotright\\Eventos e inscrições}
  \nouml{externos}{(0,-5.0)}{\guillemotleft node\guillemotright\\Ambientes externos}
    {\guillemotleft artifact\guillemotright\\Catálogo e identidade}
  \draw (cliente.east) -- node[uml/rotulo,above]{HTTPS} (servidor.west);
  \draw (servidor.south) -- node[uml/rotulo,right,text width=30mm]
    {acesso ao\\armazenamento} (armazenamento.north);
  \draw (externos.north east) -- node[uml/rotulo,above,sloped]
    {INT-001 / INT-002} (servidor.south west);
\end{tikzpicture}
}

O desenho não fixa provedor de nuvem, banco de dados ou número de máquinas.
Ambiente, recursos, dados e carga de \reqref{RNF-001} delimitam seu teste;
procedimentos de backup, recuperação e suporte são tratados no
\autoref{ap:complementar}.

\clearpage
\section{Pacotes e direção das dependências}
Na organização didática da \autoref{fig:uml-pacotes}, apresentação usa os casos
de uso; aplicação usa os conceitos e contratos do domínio; adaptadores usam
esses contratos para tratar persistência e integrações.

\diagramasrs{UML de pacotes em uma organização de código possível}{fig:uml-pacotes}
  {\begin{tikzpicture}
  \pacoteuml{apresentacao}{(0,0)}{Apresentação}
  \pacoteuml{aplicacao}{(7.5,0)}{Aplicação}
  \pacoteuml{adaptadores}{(0,-4.0)}{Adaptadores}
  \pacoteuml{dominio}{(7.5,-4.0)}{Domínio}
  \draw[uml/dependencia] (apresentacao) -- node[uml/rotulo,above,text width=24mm]{usa casos\\de uso} (aplicacao);
  \draw[uml/dependencia] (aplicacao) -- node[uml/rotulo,right]{usa conceitos e portas} (dominio);
  \draw[uml/dependencia] (adaptadores) -- node[uml/rotulo,above,text width=24mm]{usa contratos\\do domínio} (dominio);
\end{tikzpicture}
}

Os adaptadores são conectados na composição da aplicação. O pacote de domínio
não precisa conhecer HTTP, navegador ou detalhes do banco. Nomes de diretórios
e linguagens concretas pertencem ao design da implementação, não ao enunciado
de uma necessidade do usuário.

\section{Dados e rastreabilidade como modelos complementares}
A \autoref{fig:dados-logicos} expressa um modelo de dados lógico, usando
entidades, chaves e cardinalidade. O dicionário da \autoref{tab:dados} define
os tipos, a nulabilidade e as regras completas.

\diagramasrs{Modelo entidade-relacionamento de apoio, sem notação UML}{fig:dados-logicos}
  {\begin{tikzpicture}
  \node[uml/caixa,rectangle split,rectangle split parts=2,
    rectangle split part align={center,left},text width=42mm] (eventos) at (0,0)
    {\textbf{EVENTO}
    \nodepart{second}PK id\\nome\\inicio\\capacidade\\publicado};
  \node[uml/caixa,rectangle split,rectangle split parts=2,
    rectangle split part align={center,left},text width=46mm] (inscricoes) at (10.0,0)
    {\textbf{INSCRICAO}
    \nodepart{second}PK id\\FK eventoId\\participanteId (externo)\\situacao\\criadaEm\\canceladaEm (opcional)};
  \node[uml/decisao,text width=17mm] (relacao) at (5.0,0) {recebe};
  \draw (eventos.east) -- node[uml/rotulo,above]{1} (relacao.west);
  \draw (relacao.east) -- node[uml/rotulo,above]{0..N} (inscricoes.west);
  \node[uml/nota,text width=112mm] at (5.0,-4.2)
    {Cada inscrição referencia um evento. O titular pertence ao serviço de identidade:
     \texttt{participanteId} não exige uma tabela local de credenciais.
     \texttt{vagas} é derivado, não um novo dado independente.};
\end{tikzpicture}
}

\clearpage
O grafo da \autoref{fig:rastreio-modelos} ilustra uma relação entre objetivo,
regra, requisito, caso de uso e cenário de teste. É um exemplo de rastreabilidade,
sem introduzir um novo tipo de requisito ou substituir a matriz completa da
\autoref{tab:rastreabilidade}.

\diagramasrs{Relações entre uma obrigação do produto e sua verificação}{fig:rastreio-modelos}
  {\begin{tikzpicture}
  \node[uml/caixa,text width=30mm] (obj) at (0,0) {OBJ-002\\Inscrição válida};
  \node[uml/caixa,text width=30mm] (rf) at (5.1,0) {\reqref{RF-002}\\Confirmar inscrição};
  \node[uml/caixa,text width=30mm] (uc) at (10.2,0) {\ucref{UC-001}\\Realizar inscrição};
  \node[uml/caixa,text width=30mm] (rn) at (5.1,-3.2) {\rnref{RN-001}\\Capacidade e unicidade};
  \node[uml/caixa,text width=30mm] (ct) at (10.2,-3.2) {\ctref{CT-002}\\Última vaga};
  \draw[uml/fluxo] (obj) -- node[uml/rotulo,above]{motiva} (rf);
  \draw[uml/fluxo] (rf) -- node[uml/rotulo,above,text width=17mm]{detalhado\\em} (uc);
  \draw[uml/fluxo] (rn) -- node[uml/rotulo,left]{restringe} (rf);
  \draw[uml/fluxo] (ct) -- node[uml/rotulo,above]{verifica} (rn);
  \draw[uml/fluxo] (ct.north west) -- node[uml/rotulo,above,sloped]{verifica} (rf.south east);
\end{tikzpicture}
}

Na adaptação, conferir que nomes, campos, guardas e multiplicidades correspondem
às fichas. Um diagrama deve tornar uma decisão mais clara; um desenho que
contradiz o contrato precisa ser corrigido junto do texto correspondente.
